\subsection{Even Radicals} 
Even power functions are \emph{not} one-to-one, but we would like a way to reverse them, so we define a \makeblank[1in] inverse, as follows.
\begin{definition}[Radical Function (even index)]
For a positive even integer $p\geq 2$, the \term{$p$th root function}
\begin{center}$f(x)=\sqrt[p]{x}$\end{center}
is defined to be the function such that 
\begin{center}$f\inv(x)=x^p,\ x\geq 0$\end{center}
In other words, $\sqrt[p]{x}$ is the unique \emph{nonnegative} real number such that $\left(\sqrt[p]{x}\right)^p=x$.
\end{definition}
\begin{note} If we want the \emph{negative} number that, raised to the $p$th power, gives us $x$, we can simply write $-\sqrt[p]{x}$.\end{note}
\noi The domain and range of even roots is more interesting.
\begin{fact}[Domain and Range of a Radical Function (even index)]
If $f(x)=\sqrt[p]{x}$ and $f\inv(x)=x^p,\ x\geq 0$, we know that:
\begin{itemize}
\item $\dom f\inv=\makeblank[1in]$, so $\rng f=\makeblank[1in]$
\item $\rng f\inv=\makeblank[1in]$, so $\dom f=\makeblank[1in]$
\end{itemize}
\end{fact}
\begin{example}
Find the domain and range of the radical function.
\begin{lpic}{}{5}
\regtextleft{0.25}{-1}{$f(x)=3\sqrt[6]{6-x}+1$};
\foreach \x/\dr in {0/\dom,1/\rng}
	\regtextleft{\x*\value{cols}}{-5}{$\dr f={}$};
\end{lpic}
\end{example}
\begin{example}
Find each root by rewriting in terms of a power and using trial and error.
\begin{lpic}{}{3}
\foreach \y/\l/\p/\x in {1/a/4/16,2/b/2/25,3/c/4/-81}{
	\regtextright{1}{-\y}{\l. };
	\regtextleft{1}{-\y}{$\sqrt[\p]{\x}$};
}
\end{lpic}
\end{example}
Like odd roots, even roots are frequently \makeblank[1.25in].
%\begin{example}
%Use trial and error to determine what two consecutive integers each root is between.
%\begin{lpic}{}{3}
%\foreach \y/\l/\p/\x in {1/a/4/45,2/b/5/18,3/c/6/1875}{
	%\regtextright{1}{-\y}{\l. };
	%\regtextleft{1}{-\y}{$\sqrt[\p]{\x}$};
	%\regtext{2*\value{cols}-3.5}{-\y}{${}<\sqrt[\p]{\x}<{}$};
%}
%\end{lpic}
%\end{example}

